先沿着一个具体问题往下看。手里有很多图像，可以逐个读取它们，
却没有一个可直接求值的真实图像密度函数。现在希望从随机数出发，生成同类的新图像。
直接规定随机数与图像之间的复杂映射很难，扩散模型选择先做一件容易的事：
逐步给真实图像添加已知的高斯噪声，观察图像如何变得难以辨认。
然后用这些可以自行构造的训练对，学习从带噪状态回到较干净状态的方法。

这里有三个相接的问题。第一，加噪后的分布怎样写，任意时刻的样本怎样直接产生？
第二，真实的逆向分布未知，凭什么能用已知加噪过程训练它？
第三，模型训练完后，怎样把一次次局部预测连成完整生成过程？
后文先准备贝叶斯、高斯计算和随机过程的工具，再按 DDPM、NCSN、连续 SDE 的顺序回答。
DDPM 从潜变量模型与变分界出发；NCSN 从密度的空间梯度与 Langevin 采样出发。
它们在适当的噪声路径与参数换算下能够联系起来，但这个联系需要逐步推导，
不能只因为最后都在“去噪”，就把目标函数、采样算法和生成轨迹都视为相同。
